\section{模型假设}

为统一小时级调度与分钟级任务计量，假设任务开工时刻按小时粒度选择，而任务在各小时内的 GPU 和功率占用依据实际重叠时长连续折算。附件所给 GPU 需求、预计持续时间、到达时刻、来源区域、最大允许时延和最晚完成时点均被视为准确且确定的输入参数，不额外建立持续时间或需求不确定性模型。每项任务只在一个区域连续执行一次，不允许拆分、抢占、暂停或中途迁移；实时推理任务在到达时刻立即开工，弹性任务只能在到达之后且不晚于有效截止时点完成，任何任务均不得占用第 $2406$ 小时。

假设同一任务类型的单位等效 GPU 平均 IT 功率在整个研究期内保持为功率映射表给定值。非 AI IT 负荷在区域和时刻上固定且不可迁移，设施负荷由 AI 与非 AI 合计 IT 负荷乘区域 PUE 得到。区域 PUE 按附件取值，不考虑其随环境和负载率动态变化。区域间通信采用附件给出的有向单向时延，同区域时延也以附件记录为准；网络拥塞、带宽、迁移数据量、通信能耗和迁移费用均不纳入模型。

假设购电价格、售电价格、碳强度、可用新能源及固定负荷在单个小时内保持不变。各区域分别满足逐时能源平衡，不建立区域间线路潮流模型，区域售电被视为向系统边界外送。新能源可被用于直接供给负荷、储能充电、外送或弃电，其分配由约束与评价目标共同决定。运行碳排放仅按电网购电量与对应碳强度计算，不计新能源、储能设备及通信设施的生命周期排放。

储能初始 SOC 被固定为储能信息表给出的第 $0$ 小时运行前状态，逐时基准 SOC 仅用于结果对比。充电效率与放电效率分别作用于 SOC 增量和减量，并忽略题面未给出的自放电。第 $2406$ 小时运行后的 SOC 不得低于初始 SOC。同一区域同一小时不得同时充电和放电；购电与售电均受给定边界限制，售电仅来自当期新能源分流，不允许通过购电转售形成套利。

问题三中的 AI IT 负荷与非 AI IT 负荷被视为固定输入，不重新调整任务执行区域和开工时刻。问题一测试区间的预测结果仅用于预测性能评价，末端基础调度以及问题二至问题四的正式计算均采用实际任务。服务质量由任务即时开工、网络时延 SLA、按期完成和等待时间等可观测信息构造，不被解释为真实用户满意度。情景分析仅改变指定研究因素，其余任务、设备、网络、功率映射和约束保持一致。普通数据空白不被自动补零，只有在字段互斥关系、备注及相邻记录共同支持时，才被判定为结构零。
